\documentclass[stu,12pt]{apa7}

\usepackage[spanish,mexico]{babel}
\usepackage[style=apa]{biblatex}
\usepackage{geometry}
\usepackage{bookmark}
\usepackage{csquotes}
\usepackage{graphicx}
\usepackage{hyperref}
\usepackage{listings}
\usepackage{xcolor}
\usepackage{import}
\usepackage{xifthen}
\usepackage{pdfpages}
\usepackage{multirow}
\usepackage{transparent}
\usepackage{mathptmx}
\usepackage{amsmath}
\usepackage{amssymb}

\addbibresource{referencias.bib}
\affiliation{Facultad de Inform\'atica, Universidad Aut\'onoma de Quer\'etaro}
\course{G31, Semestre 2, 2014: Desarrollo Humano II}
\professor{Servin González Cecilia}

\hypersetup{
	pdftitle={Nombre de Proyecto Final},
	colorlinks=false,
}
\title{El impacto que genera la Inteligencia M\'ultiple Intrapersonal y su
Conducta durante el Desarrollo Profesional}
\shorttitle{Inteligencia M\'ultiple, Conducta en el Desarrollo Personal}
\author{Araiza Cort\'ez Jos\'e Antonio, Baqueiro Lezama Mateo, Bautista Olvera
Marco El\'ias, C\'ardenas Camacho Alexis Emilio, Garc\'ia Mej\'ia Miguel
\'Angel, Gonz\'alez C\'ordoba Bryan Michel, Ju\'arez Ram\'irez Gabriel, Le\'on
L\'opez Michel Christopher, Mac\'ias Fonseca Alejandro, Mata Guerra David, 
Mondrag\'on Fern\'andez Josu\'e Oziel, N\'u\~nez Rodr\'iguez C\'esar de Jes\'us,
Olvera Garc\'ia Marco Antonio, Padr\'on Moya Luis Gabriel, Pichardo Botillo
Blanca Isabel, Pi\~n\'on Garc\'ia Paola, Rodr\'iguez Mart\'inez David Emmanuel,
Serv\'in L\'opez Emilio, Ugalde Arias Alison, Velasco Rom\'an Paola Esmeralda,
Zamorano Monzalvo Daniel}
\duedate{16 de mayo de 2024\\
	\begin{center}
	\includegraphics[width=2cm]{./assets/img/escudo-uaq.png}\end{center}
	\begin{center}\includegraphics[width=2cm]
	{./assets/img/escudo-uaq-facultad-informatica.png}
	\end{center}
}

\begin{document}
\maketitle
\tableofcontents
\section{Introducci\'on}
\input{./secciones/03-introduccion.tex}
\section{Planteamiento del Problema}
\input{./secciones/04-planteamiento-del-problema.tex}
\section{Objetivo General y Particular}
\input{./secciones/05-objetivo-general-y-particular.tex}
\section{Justificaci\'on}
\documentclass{article}
\begin{document}

\section{Justificación}
Este artículo surge por dos motivos fundamentales. En primer lugar, busca recrear un ambiente laboral en el aula con el objetivo de promover la experiencia del trabajo en equipo. Nos interesa observar cómo el compromiso, la responsabilidad, la gestión del tiempo y el esfuerzo conjunto pueden llevarnos a alcanzar objetivos más ambiciosos que cuando trabajamos de manera individual o en equipos desorganizados.

En segundo lugar, este estudio pretende resaltar la importancia de la introspección en el entorno laboral. A menudo se enfatiza la relevancia de habilidades interpersonales como la comunicación y la expresión, pero rara vez se aborda el valor de la introspección. Este artículo busca demostrar el impacto significativo de la inteligencia intrapersonal en el desarrollo de las''soft skill'' y su influencia positiva en el desempeño laboral.

Al explorar la interacción entre el trabajo en equipo y la introspección, este artículo busca ofrecer una perspectiva más completa sobre cómo las habilidades intrapersonales pueden complementar y potenciar las habilidades interpersonales en un entorno laboral dinámico.

\end{document}


\section{Hip\'otesis}
\input{./secciones/07-hipotesis.tex}
\section{Marco Te\'orico}
\input{./secciones/08-marco-teorico.tex}
\section{Metodolog\'ia de la Investigaci\'on}
\input{./secciones/09-metodologia-de-la-investigacion.tex}
\section{Poblaci\'on y Muestra}
\input{./secciones/10-poblacion-y-muestra.tex}
\section{T\'ecnicas e Instrumentos de Investigaci\'on}
\input{./secciones/11-tecnicas-e-instrumentos-de-investigacion.tex}
\section{Recopilaci\'on de Datos}
\input{./secciones/12-recopilacion-de-datos}
\section{An\'alisis de Resultados}
\input{./secciones/13-analisis-de-resultados}
\section{Presentaci-'on de Resultados y Conclusi\'on}
\input{./secciones/14-presentacion-de-resultados-y-conclusion.tex}
\begin{appendix}
\input{./secciones/15-anexos}
\end{appendix}
\printbibliography
\end{document}
